\section{Deferred statements}
\label{supp:statements}

This section states the results that the main text uses through pointers.
Numbering is internal to the Supplement; references such as
Theorem~\ref{thm:rep} point to the main text.



\begin{proposition}[Cycle densities are power sums]
\label{sprop:ident}
Let $X\sim P^{\mathrm{net}}_{\sigma,\blam}$ with $\sum_k\lam_k^2<\infty$.
For distinct indices $i_1,\dots,i_m$, $m\ge3$,
\begin{equation*}
\E\bigl[X_{i_1i_2}X_{i_2i_3}\cdots X_{i_mi_1}\bigr] \;=\; p_m(\blam)
\;=\;\sum_k \lam_k^m ,
\qquad
\E\bigl[X_{12}^2\bigr] \;=\; \sigma^2 + p_2(\blam).
\end{equation*}
More generally, for a finite simple graph $F$ with at least one edge:
$\E\prod_{e\in F}X_e=0$ whenever some edge of $F$ lies in no cycle of $F$
(in particular for every tree), and
$\E\prod_{e\in F}X_e=\prod_{j\le r}p_{|C_j|}(\blam)$ when $F$ is a
vertex-disjoint union of cycles $C_1,\dots,C_r$, each of length at least
three. The product formula does not extend further: two triangles sharing
one vertex have expectation $p_3(\blam)^2+2\,p_6(\blam)$, and a doubled
edge has expectation $\sigma^2+p_2(\blam)$. Consequently the noise level
and the multiset of nonzero eigenvalues, $(\sigma^2,\{\lam_k\})$, are
identified by the ergodic law, and two parameter pairs yield the same
network law if and only if they agree in this sense.
\end{proposition}

\begin{proposition}[Emergence and estimation]
\label{sprop:bbp}
Let $\sigma>0$, let $X\sim R_{0,\sigma,\blam}$ have finitely many
nonzero $\lam_1\ge\dots\ge\lam_K$ (all $\ne0$), and let
$\mu_{1}\ge\dots\ge\mu_n$ be the eigenvalues of the fully observed
corner $X_n/(\sigma\sqrt n)$; the zero-diagonal observation changes
every $\mu_j$ by $O_P(\log n/\sqrt n)$ and all statements below are
robust to it. Set $s_k=\sqrt n\,\lam_k/\sigma$. Then, as $n\to\infty$
with the array fixed:
\begin{enumerate}[label=(\roman*),itemsep=1pt,topsep=2pt]
\item (Resolution.) For every $k$ with $\lam_k>0$,
$\mu_k/s_k\toas1$ and
$\mu_k-(s_k+1/s_k)=s_k\bigl(\lVert\Xi_k\rVert^2/n-1\bigr)+o_P(1)
=O_P(1)$, the $O_P(1)$ term reflecting the realized factor norm; the
mirrored statements hold for $\lam_k<0$ at the smallest eigenvalues.
For every fixed $\varepsilon>0$, the number of eigenvalues outside
$[-2-\varepsilon,2+\varepsilon]$ converges a.s.\ to
$\#\{k:\lam_k\neq0\}$.
\item (Finite-$n$ boundary.) Against local sequences
$\lam^{(n)}=c\,n^{-1/2}$ the transition is at $c=\sigma$: for
$c<\sigma$ no eigenvalue separates from the bulk and no consistent
test exists along the contiguous sequences
\citep{onatski2013asymptotic,perry2018optimality}; for $c>\sigma$ the
top eigenvalue separates and its eigenvector correlates with the
planted factor, with squared overlap $\to 1-c^{-2}\sigma^2$.
\item (Plug-in estimation.) The empirical spectral distribution of
$X_n/\sqrt n$ converges weakly a.s.\ to the semicircle law of radius
$2\sigma$, so
$\hat\sigma=\hat q_{3/4}(\spec(X_n/\sqrt n))/q_{3/4}^{\rm sc}$,
$q_{3/4}^{\rm sc}\approx0.8079$, satisfies $\hat\sigma\toas\sigma$.
For each $j$ with $|\mu_j|>2+\varepsilon_n$ set
$\hat\lam_j=\operatorname{sgn}(\mu_j)\,
\tfrac{\hat\sigma}{\sqrt n}\,\hat s_j$ with
$\hat s_j=\bigl(|\mu_j|+\sqrt{\mu_j^2-4}\bigr)/2$; then
$\hat\lam_k/\lam_k\toas1$ for every $k$ with $\lam_k\ne0$, and, with
$\varepsilon_n\downarrow0$ satisfying
$n^{2/3}\varepsilon_n\to\infty$ and
$\varepsilon_n\sqrt n/\log n\to\infty$,
$\hat K=\#\{j:|\mu_j|>2+\varepsilon_n\}$ is a consistent estimator of
the number of spikes.
\end{enumerate}
\end{proposition}

\begin{corollary}[Non-spectral summaries are conditionally parameter-free]
\label{scor:ancillary}
Let $S:\Sym_n\to\R^q$ be any measurable network statistic. Under any
family $\mathcal{P}$ of orthogonally invariant laws:
(i) the conditional law of $S(X)$ given $D(X)$ is parameter-free and
exactly simulable by Haar rotation of the observed spectrum;
(ii) $\E[S(X)\mid D]$ is a spectral statistic, and for any parameter
$\theta$ of any subfamily, any estimator or test based on $S$ is weakly
dominated by one based on $D$ alone (Rao--Blackwell; convex losses,
tests via the conditional test function);
(iii) on the simple-spectrum event, any statistic of the frame $V(X)$
alone is ancillary: $V(X)\sim\nu_n$, the sign-normalized Haar law of
Theorem~\ref{thm:suff}(iii), whose distribution is free of the
parameter.
\end{corollary}

\begin{proposition}[Consistency against localized spikes]
\label{sprop:consistency}
Let $X^{(n)}$ follow the spiked model
$X^{(n)}=\tfrac{\sigma}{\sqrt n}W_n+\theta v_nv_n^\top$ with $\theta>\sigma$
and unit vectors $v_n$ satisfying
$\liminf_n \lVert v_n\rVert_4^4 \ge c_0>0$
(e.g.\ $v_n$ supported on $O(1)$ many coordinates, or any hub-like
direction). Let $p^{(n)}$ be the conditional $p$-value of
$T_{\mathrm{IPR}}$ with $B_n\to\infty$. Then $p^{(n)}\toP0$; the test is
consistent. Under the rotatable null with the same spectrum,
$T_{\mathrm{IPR}}=O_P(\log n/n)$, so the separation is by a factor of
order $n/\log n$.
\end{proposition}

\begin{corollary}[Spherical spikes are null; localized detection threshold]
\label{scor:boundary}
Fix $\sigma>0$.
(i) If $v$ is uniform on the unit sphere, independent of $W_n$, then
$\sigma W_n/\sqrt n+\theta vv^\top$ is orthogonally invariant for
every $\theta$, hence lies in the composite null of
Theorem~\ref{thm:test}; the conditional test has exact level against
it at every spike strength.
(ii) For $\ell_4$-nondegenerate directions
($\liminf\lVert v_n\rVert_4^4\ge c_0>0$) and every fixed
$s=\theta/\sigma>1$, the IPR test is consistent
(Proposition~\ref{sprop:consistency}).
(iii) For every fixed $s<1$ with spherical direction, the alternatives
are contiguous to the point null with $\sigma$ fixed and $\blam=0$
\citep{perry2018optimality}; every test sequence whose size tends to
zero has power tending to zero, so no consistent test exists.
Sparse priors can admit sub-boundary consistent detection by
exhaustive scans outside this battery, and positive-size tests below
the boundary may retain nontrivial power, which contiguity does not
exclude.
\end{corollary}

\begin{remark}[Hollow laws and Haar orbits are mutually singular]
\label{srem:hollowtv}
The requirement that the diagonal be retained is essential. Suppose
that \(Y\in\Sym_n\) is invariant under orthogonal conjugation and is
hollow almost surely. For every deterministic unit vector
\(a\in\mathbb R^n\), choose \(q_a\in\OO(n)\) such that
\(q_a^\top e_1=a\). Orthogonal invariance gives
\[
a^\top Ya
=
e_1^\top q_aYq_a^\top e_1
\eqd
e_1^\top Ye_1
=
0.
\]
Let \(\mathcal A\) be a countable dense subset of the unit sphere.
With probability one,
\[
a^\top Ya=0
\qquad\text{for every }a\in\mathcal A.
\]
For each realization in this probability-one event, continuity of
\(a\mapsto a^\top Ya\) extends the equality to the entire unit sphere,
and homogeneity extends it to all \(a\in\mathbb R^n\). Since a real
symmetric matrix is determined by its quadratic form, \(Y=0\) almost
surely.
Thus no nondegenerate hollow law, including a nondegenerate
zero-completed network law
\(P^{\mathrm{net}}_{\sigma,\blam}\), can be exactly invariant under
all orthogonal conjugations. Theorem~\ref{thm:suff} applies exactly to
the corresponding fully observed symmetric corner before hollowing.
There is also a stronger obstruction in total variation. Let
\[
\mathcal H_n
:=
\left\{
x\in\Sym_n:
x_{11}=\cdots=x_{nn}=0
\right\}.
\]
For every \(d\in\Delta_n\setminus\{0\}\),
\[
\kappa(d,\mathcal H_n)=0.
\]
For \(n=1\), this is immediate. For \(n\geq2\), if
\(U\sim\Haar_n\), then its first row is uniformly distributed on the
unit sphere and
\[
\left(U\diag(d)U^\top\right)_{11}
=
\sum_{i=1}^n d_iU_{1i}^2.
\]
As a function of the first row of \(U\), this is a nonzero
real-analytic function on the unit sphere. Its zero set has Haar
surface measure zero. Hence the Haar orbit belongs to
\(\mathcal H_n\) with probability zero.
In contrast, suppose that \(X\) is hollow almost surely and let
\(Q(d,\cdot)\) be a regular conditional law of \(X\) given
\(D(X)=d\). Then
\[
Q(d,\mathcal H_n)=1
\qquad
\text{for \(P_D\)-almost every }d.
\]
Under the convention
\[
d_{\mathrm{TV}}(\mu,\eta)
:=
\sup_{A\in\mathcal B(\Sym_n)}
|\mu(A)-\eta(A)|,
\]
it follows that
\[
d_{\mathrm{TV}}\!\left(
Q(d,\cdot),\kappa(d,\cdot)
\right)
=
1
\qquad
\text{for \(P_D\)-almost every }d\neq0.
\]
Therefore an exactly hollow nonzero observation cannot satisfy a
nontrivial conditional total-variation approximation to the Haar
orbit. Any useful approximation must use a weaker metric, a diagonal
augmentation model, or a different observation model.
\end{remark}

\begin{remark}[Fixed-diagonal observations]
\label{srem:fixeddiag}
The conclusion is nontrivial only when \(\varepsilon<1\). Fix
\(c\in\mathbb R\), and define
\[
\mathcal H_{n,c}
:=
\left\{
x\in\Sym_n:
x_{11}=\cdots=x_{nn}=c
\right\}.
\]
If \(P(\mathcal H_{n,c})=1\), then
\[
P_d(\mathcal H_{n,c})=1
\qquad
\text{for \(P_D\)-almost every \(d\)}.
\]
On the other hand, if \(d\neq c\mathbf 1_n\), then
\[
\kappa_d(\mathcal H_{n,c})=0.
\]
Indeed, if \(V\) is the first row of a Haar-distributed orthogonal
matrix, then \(V\) is uniform on the unit sphere and
\[
\left(U\diag(d)U^\top\right)_{11}
=
\sum_{i=1}^{n}d_iV_i^2.
\]
On the unit sphere, the equality
\[
\sum_{i=1}^{n}d_iV_i^2=c
\]
is equivalent to
\[
\sum_{i=1}^{n}(d_i-c)V_i^2=0.
\]
If the restriction of this polynomial to the sphere were identically
zero, evaluation at the coordinate vectors would give
\(d_i=c\) for every \(i\). Thus, when \(d\neq c\mathbf 1_n\), it is a
nonzero real-analytic function on the sphere. For \(n\geq2\), its zero
set has surface measure zero; the case \(n=1\) is immediate.
Therefore
\[
d_{\mathrm{TV}}(P_d,\kappa_d)=1
\qquad
\text{for \(P_D\)-almost every }
d\in\Delta_n\setminus\{c\mathbf 1_n\}.
\]
Consequently, if \(P(\mathcal H_{n,c})=1\) and
\eqref{eq:approx-haar-conditionals} holds with \(\varepsilon<1\), then
\[
P_D\{c\mathbf 1_n\}=1.
\]
Every real symmetric matrix whose eigenvalues all equal \(c\) is
\(cI_n\), so
\[
P\{cI_n\}=1.
\]
Thus the full-TV assumption cannot yield a nontrivial guarantee for
an intrinsically hollow nonzero matrix or for a nonidentity
correlation matrix with fixed unit diagonal. Common preprocessing is
permitted only after beginning with an unprocessed matrix satisfying
\eqref{eq:approx-haar-conditionals}; preprocessing cannot itself make
an intrinsically fixed-diagonal observation satisfy that assumption.
\end{remark}
